\documentclass[11pt]{article}

\usepackage[margin=1in]{geometry}
\usepackage{amsmath,amssymb,amsthm,mathtools}
\usepackage[T1]{fontenc}
\usepackage{booktabs}
\usepackage{xcolor}
\usepackage[colorlinks=true,linkcolor=blue,urlcolor=blue]{hyperref}

\newtheorem{theorem}{Theorem}[section]
\newtheorem{corollary}[theorem]{Corollary}
\newtheorem{proposition}[theorem]{Proposition}
\theoremstyle{definition}
\newtheorem{definition}[theorem]{Definition}
\theoremstyle{remark}
\newtheorem{remark}[theorem]{Remark}

% status macros
\newcommand{\leanname}[1]{\texttt{#1}}
\newcommand{\proved}{\textcolor{green!55!black}{\textsf{proved, axiom-clean}}}
\newcommand{\cited}{\textcolor{orange!85!black}{\textsf{cited axiom}}}
\newcommand{\lstatus}[2]{\par\smallskip\noindent{\footnotesize $\vartriangleright$ \emph{Lean:} \leanname{#1} \;---\; #2.}}

% notation
\newcommand{\R}{\mathbb{R}}
\newcommand{\xij}{\xi^j}
\newcommand{\zj}{\zeta^j}
\newcommand{\Iint}{I}
\DeclareMathOperator*{\argmin}{arg\,min}

\title{\textbf{Inverse Green's Learning}\\[2pt]
  \large A Lean 4 / Mathlib formalization of the factorization theory}
\author{Formalization notes (IGL project)}
\date{\today}

\begin{document}
\maketitle

\begin{abstract}
We record the machine-checked content of the \emph{Inverse / Intrinsic Green's Learning}
(IGL) factorization theory, formalized in Lean~4 against Mathlib (toolchain
\leanname{leanprover/lean4:v4.30.0}). The factorization theory is proved and kernel-checked:
the rank-1 and $d$-fold separable factorizations; the general \emph{Master Formula} with its
Additive-Green and Kolmogorov--Arnold special cases; gauge invariance under
measure-preserving reparametrization; the \emph{spectral} factorized integral; the
\emph{CP-rank} sub-additivity and sub-multiplicativity bounds with rank-friendly
differentiation; and the \emph{separability barrier} (sufficiency, with a concrete necessity
witness) and its compensation mechanism. Each carries its Lean name and its complete axiom
dependency; the proved results rest only on Mathlib's three standard
axioms (\leanname{propext}, \leanname{Classical.choice}, \leanname{Quot.sound}). The
classical elliptic-PDE prerequisites absent from Mathlib --- the Green's kernel itself and
the Braess--Hackbusch exponential-sum rank bound --- are introduced as \emph{named, cited
axioms}, so every result that depends on them names them explicitly.
\end{abstract}

\section{Setup}

Fix an intrinsic dimension $d$ and ranks $K, R \in \mathbb{N}$. Points of the coordinate
space are $\xi = (\xi^1,\dots,\xi^d) \in \R^d$ (the query / intrinsic coordinate) and
$\zeta = (\zeta^1,\dots,\zeta^d) \in \R^d$ (the integration variable). We integrate against
Lebesgue measure on $\R^d$, which is the product $\bigotimes_{j=1}^d \mathrm{vol}$ of the
one-dimensional Lebesgue measures (in Mathlib, \leanname{MeasureTheory.volume\_pi}, a
definitional equality).

\begin{definition}[Rank-$R$ separable source]\label{def:source}
$\displaystyle \hat f(\zeta) = \sum_{r=1}^{R} w_r \prod_{j=1}^{d} \varphi_{r,j}(\zj)$,
with weights $w_r \in \R$ and univariate factors $\varphi_{r,j}\colon \R \to \R$.
\lstatus{IGL.separableSource}{definition}
\end{definition}

\begin{definition}[Rank-$K$ separable kernel]\label{def:kernel}
$\displaystyle G(\xi,\zeta) = \sum_{k=1}^{K} \gamma_k \prod_{j=1}^{d} G_{k,j}(\xij,\zj)$,
with mixing weights $\gamma_k \in \R$ and univariate kernels $G_{k,j}\colon \R\times\R\to\R$.
\lstatus{IGL.separableKernel}{definition}
\end{definition}

\begin{definition}[One-dimensional partial integral]\label{def:I}
$\displaystyle \Iint_{k,r,j}(\xij) = \int_{\R} G_{k,j}(\xij,\zj)\,\varphi_{r,j}(\zj)\,d\zj$.
\lstatus{IGL.partialIntegral}{definition}
\end{definition}

\begin{definition}[IGL output]\label{def:output}
$\displaystyle u(\xi) = \sum_{k=1}^{K}\sum_{r=1}^{R} \gamma_k w_r \prod_{j=1}^{d} \Iint_{k,r,j}(\xij)$.
\lstatus{IGL.iglParticular}{definition}
\end{definition}

\section{The factorization theorems}

\begin{theorem}[Rank-1 factorization]\label{thm:rank1}
For $f,g\colon\R\to\R$,
$\displaystyle \int_{\R^2} f(p_1)\,g(p_2)\,d(p_1,p_2) = \Big(\int_\R f\Big)\Big(\int_\R g\Big).$
\end{theorem}
\begin{proof}
Immediate from the product-measure integral of a separated product; in Mathlib this is
\leanname{MeasureTheory.integral\_prod\_mul}, which holds unconditionally (both sides
vanish when the factors fail to be integrable).
\end{proof}
\lstatus{igl\_separable\_factorization}{\proved}

\begin{theorem}[$d$-fold separable factorization]\label{thm:dfold}
For $h_1,\dots,h_d\colon\R\to\R$,
\[
  \int_{\R^d} \prod_{j=1}^{d} h_j(\zj)\,d\zeta = \prod_{j=1}^{d}\int_\R h_j(\zj)\,d\zj .
\]
\end{theorem}
\begin{proof}
This is the Fubini--Tonelli factorization for a product measure: since Lebesgue measure on
$\R^d$ equals $\bigotimes_j \mathrm{vol}$ and each factor is $\sigma$-finite, the integral
of a product of coordinate-wise factors equals the product of the one-dimensional
integrals. In Mathlib this is \leanname{MeasureTheory.integral\_fintype\_prod\_eq\_prod}.
Generalizes Theorem~\ref{thm:rank1} from two factors to $d$.
\end{proof}
\lstatus{IGL.fubini\_dfold}{\proved}

\begin{theorem}[Master Formula]\label{thm:master}
Let $\hat f$ and $G$ be separable as in Definitions~\textup{\ref{def:source}} and
\textup{\ref{def:kernel}}. Assume that for each $k,r$ the function
$\zeta \mapsto \big(\prod_j G_{k,j}(\xij,\zj)\big)\big(\prod_j \varphi_{r,j}(\zj)\big)$ is
Lebesgue-integrable on $\R^d$. Then
\[
  \int_{\R^d} G(\xi,\zeta)\,\hat f(\zeta)\,d\zeta
  = \sum_{k=1}^{K}\sum_{r=1}^{R} \gamma_k w_r \prod_{j=1}^{d} \Iint_{k,r,j}(\xij)
  = u(\xi),
\]
reducing the per-sample cost from $O(N^d)$ to $O(KRd)$.
\end{theorem}
\begin{proof}
Write the integrand with the separable forms of kernel and source:
\[
  \int_{\R^d}
    \Big(\sum_{k}\gamma_k\prod_j G_{k,j}(\xij,\zj)\Big)
    \Big(\sum_{r} w_r\prod_j \varphi_{r,j}(\zj)\Big)d\zeta .
\]

\emph{Step 1 (expand the product of sums).} Distributing the two finite sums turns the
integrand into $\sum_{k}\sum_{r}\big(\gamma_k\prod_j G_{k,j}\big)\big(w_r\prod_j \varphi_{r,j}\big)$.

\emph{Step 2 (interchange the finite sums and the integral).} Each summand equals
$\gamma_k w_r \cdot \big(\prod_j G_{k,j}\big)\big(\prod_j \varphi_{r,j}\big)$, i.e.\ the
hypothesized integrand scaled by the constant $\gamma_k w_r$, hence integrable. As a finite
sum of integrable functions is integrable and the Bochner integral is linear,
\[
  \int_{\R^d}\sum_{k,r}(\cdots)\,d\zeta
  = \sum_{k}\sum_{r}\int_{\R^d}
      \Big(\gamma_k\prod_j G_{k,j}\Big)\Big(w_r\prod_j \varphi_{r,j}\Big)d\zeta .
\]
This is the \emph{only} place the integrability hypothesis is used.

\emph{Step 3 (factor each $(k,r)$ term).} Fix $k,r$. Pull out the scalar $\gamma_k w_r$ and
merge the two $d$-fold products into one product of $d$ univariate factors:
\[
  \int_{\R^d}\Big(\gamma_k\prod_j G_{k,j}\Big)\Big(w_r\prod_j \varphi_{r,j}\Big)d\zeta
  = \gamma_k w_r \int_{\R^d}\prod_{j=1}^{d}\big(G_{k,j}(\xij,\zj)\,\varphi_{r,j}(\zj)\big)\,d\zeta .
\]
The integrand is now a product of $d$ factors, the $j$-th depending only on $\zj$, so
Theorem~\ref{thm:dfold} (with $h_j(\zj) = G_{k,j}(\xij,\zj)\varphi_{r,j}(\zj)$) gives
$\gamma_k w_r\prod_j \int_\R G_{k,j}(\xij,\zj)\varphi_{r,j}(\zj)\,d\zj
 = \gamma_k w_r\prod_j \Iint_{k,r,j}(\xij)$,
using Definition~\ref{def:I}. Summing over $k,r$ yields the claim.
\end{proof}
\lstatus{IGL.fubini\_factorization}{\proved}

\section{Special cases}

\begin{corollary}[Additive-Green, $K=1$]\label{cor:additive}
For a single separable kernel $\prod_j G_j(\xij,\zj)$ and a rank-$R$ separable source (under
the analogous integrability hypothesis),
$\displaystyle \int_{\R^d}\Big(\prod_j G_j(\xij,\zj)\Big)\hat f(\zeta)\,d\zeta
  = \sum_{r} w_r \prod_{j=1}^{d} \Iint_{r,j}(\xij).$
\end{corollary}
\begin{proof}
The $K=1$ instance of Theorem~\ref{thm:master}: there is no $k$-sum, so Steps~1--2 reduce to
pulling the single finite $r$-sum through the integral (linearity, justified by
integrability), and Step~3 applies Theorem~\ref{thm:dfold} to each $r$-term verbatim.
\end{proof}
\lstatus{IGL.additiveGreen\_factorization}{\proved}

\begin{corollary}[Kolmogorov--Arnold limit, $G=\mathrm{Id}$]\label{cor:kan}
If each one-dimensional integral reproduces its source factor,
$\Iint_{k,r,j}(\xij) = \varphi_{r,j}(\xij)$ (the identity-/reproducing-kernel condition),
then the output collapses to a sum of products of univariate functions of the coordinates:
$\displaystyle u(\xi) = \sum_{k}\sum_{r} \gamma_k w_r \prod_{j=1}^{d}\varphi_{r,j}(\xij).$
\end{corollary}
\begin{proof}
Substitute the reproducing identity $\Iint_{k,r,j}(\xij)=\varphi_{r,j}(\xij)$ into
Definition~\ref{def:output}.
\end{proof}
\lstatus{IGL.kan\_reduction}{\proved}

\section{Gauge invariance}

\begin{theorem}[Gauge invariance]\label{thm:gauge}
Let $e\colon \R^d \to \R^d$ be a measure-preserving reparametrization (a measurable
equivalence with $e_\sharp \mathrm{vol} = \mathrm{vol}$; e.g.\ any rotation, as used by the
random-rotation embedding). Then for every $h\colon\R^d\to\R$,
$\displaystyle \int_{\R^d} h(e(\zeta))\,d\zeta = \int_{\R^d} h(\zeta)\,d\zeta.$
In particular the IGL prediction integral is invariant under such reparametrizations; the
``source transforms inversely'' of the informal account is the change of variables
$h \circ e$.
\end{theorem}
\begin{proof}
Change of variables for a measure-preserving map: in Mathlib,
\leanname{MeasureTheory.MeasurePreserving.integral\_comp'} applied to $e$ and $h$.
\end{proof}
\lstatus{IGL.gauge\_invariance}{\proved}

\section{Spectral form}

\begin{theorem}[Additive eigenvalues become a product]\label{thm:expneg}
For $\alpha\in\R$ and $\lambda_1,\dots,\lambda_d\in\R$,
$\displaystyle \exp\!\Big(-\alpha\sum_{j=1}^d\lambda_j\Big)=\prod_{j=1}^d\exp(-\alpha\lambda_j).$
\end{theorem}
\begin{proof}
Distribute the scalar over the finite sum and apply $\exp(\sum_j a_j)=\prod_j\exp(a_j)$
(Mathlib \leanname{Real.exp\_sum}) with $a_j=-\alpha\lambda_j$.
\end{proof}
\lstatus{IGL.exp\_neg\_sum\_eigs}{\proved}

\begin{theorem}[Spectral factorized integral]\label{thm:spectral}
Let modes be multi-indices $n\colon\{1,\dots,d\}\to\{1,\dots,M\}$, with per-axis eigenvalues
$\lambda_{j,m}$, eigenfunctions $\varphi_{j,m}$, and a source separable in the eigenbasis,
$\hat f_n=\prod_j g_{j,n_j}$. With the exponential-sum surrogate
$1/\lambda\approx\sum_k\gamma_k e^{-\alpha_k\lambda}$ applied modewise,
\[
  \sum_{k}\gamma_k\sum_{n}\Big(\prod_j e^{-\alpha_k\lambda_{j,n_j}}\varphi_{j,n_j}(\xij)\Big)\prod_j g_{j,n_j}
  = \sum_{k}\gamma_k\prod_{j=1}^d\sum_{m=1}^M e^{-\alpha_k\lambda_{j,m}}\varphi_{j,m}(\xij)\,g_{j,m}.
\]
\end{theorem}
\begin{proof}
Fix $k$ and merge the two products over $j$ into a single
$\prod_j\big(e^{-\alpha_k\lambda_{j,n_j}}\varphi_{j,n_j}(\xij)\,g_{j,n_j}\big)$. The sum over
multi-indices $n$ of a product of per-axis factors equals the product over axes of the
per-axis sums (the distributive law $\prod_j\sum_m=\sum_{n}\prod_j$, Mathlib
\leanname{Finset.prod\_univ\_sum}) --- the discrete analogue of Theorem~\ref{thm:dfold}. The
additive eigenvalue coupling was converted into a product beforehand by
Theorem~\ref{thm:expneg}.
\end{proof}
\lstatus{IGL.spectral\_factorization}{\proved}

\section{CP-rank and rank-friendly differentiation}

\begin{definition}[CP-rank, upper bound]\label{def:cprank}
$\mathrm{HasCPRank}(f,n)$ holds when $f(\zeta)=\sum_{i} w_i\prod_j\varphi_{i,j}(\zj)$ over an
index set of cardinality $n$ --- a \emph{constructive} upper bound on the CP-rank: it carries
an explicit $n$-term decomposition (no minimisation), which is all the paper's bounds need.
\lstatus{IGL.HasCPRank}{definition}
\end{definition}

\begin{theorem}[Sub-additivity of CP-rank]\label{thm:rankadd}
$\mathrm{rank}(f+g)\le\mathrm{rank}(f)+\mathrm{rank}(g).$
\end{theorem}
\begin{proof}
Concatenate the two term families over the disjoint union of the index sets; a sum over a
disjoint union splits into the two sums (\leanname{Fintype.sum\_sum\_type}) and cardinalities
add (\leanname{Fintype.card\_sum}).
\end{proof}
\lstatus{IGL.HasCPRank.add}{\proved}

\begin{theorem}[Sub-multiplicativity of CP-rank]\label{thm:rankmul}
$\mathrm{rank}(f\cdot g)\le\mathrm{rank}(f)\cdot\mathrm{rank}(g)$ for the pointwise product.
\end{theorem}
\begin{proof}
Index the result by pairs. Expand the product of the two CP sums
(\leanname{Finset.sum\_mul\_sum}), repackage the resulting double sum as a sum over the
product index (\leanname{Fintype.sum\_prod\_type}), and merge the two $j$-products into one
(\leanname{Finset.prod\_mul\_distrib}); cardinalities multiply (\leanname{Fintype.card\_prod}).
\end{proof}
\lstatus{IGL.HasCPRank.mul}{\proved}

\begin{theorem}[Differentiation is rank-friendly]\label{thm:diffrank}
A coordinate-wise separable operator $L=\sum_j L_j$ (each $L_j$ acting on the $j$-th factor)
maps a rank-$R$ source to one of rank at most $R\,d$:
$\displaystyle L\hat f=\sum_r w_r\sum_j (L_j\varphi_{r,j})\prod_{i\ne j}\varphi_{r,i},$
a sum of $R d$ separable terms.
\end{theorem}
\begin{proof}
The $(r,j)$ term replaces the $j$-th factor of the $r$-th source term by its $L_j$-image.
Re-indexing the $R$-sum-of-$d$-sums as a single sum over
$\{1,\dots,R\}\times\{1,\dots,d\}$ (\leanname{Fintype.sum\_prod\_type},
\leanname{Finset.mul\_sum}) exhibits the decomposition; the index set has cardinality $Rd$.
\end{proof}
\lstatus{IGL.applyCoordOp\_hasCPRank}{\proved}

\section{The separability barrier}

\begin{definition}[Metric volume factor]\label{def:metric}
On a curved chart the integral carries $\sqrt{|g(\zeta)|}$:
$u(\xi)=\int_{\R^d} G(\xi,\zeta)\,\tilde f(\zeta)\,\sqrt{|g(\zeta)|}\,d\zeta$.
\end{definition}

\begin{theorem}[Barrier --- sufficiency (``if'')]\label{thm:barrierif}
If the metric factor is itself separable, $\sqrt{|g(\zeta)|}=\prod_j\chi_j(\zj)$, then for a
rank-$R$ separable source (under the analogous integrability hypothesis) the compensated
integral still factorizes:
\[
  \int_{\R^d}\Big(\prod_j G_j(\xij,\zj)\Big)\Big(\sum_r w_r\prod_j\varphi_{r,j}(\zj)\Big)\Big(\prod_j\chi_j(\zj)\Big)d\zeta
  = \sum_r w_r\prod_{j=1}^d\int_\R G_j\,\varphi_{r,j}\,\chi_j\,d\zj.
\]
\end{theorem}
\begin{proof}
Distribute the finite $r$-sum through the integral (linearity, justified by integrability);
for each $r$ merge the three products into $\prod_j\big(G_j\varphi_{r,j}\chi_j\big)$
(\leanname{Finset.prod\_mul\_distrib}) and apply Theorem~\ref{thm:dfold}. The metric is
absorbed axis-by-axis: a separable $\sqrt{|g|}$ does not obstruct the collapse.
\end{proof}
\lstatus{IGL.barrier\_if\_factorization}{\proved}

\begin{theorem}[Barrier --- necessity (witness)]\label{thm:barrieronly}
The metric factor $\sqrt{|g(\zeta)|}=1+\zeta_0\zeta_1$ on $\R^2$ is \emph{not} rank-1: there
is no decomposition $1+\zeta_0\zeta_1=a(\zeta_0)\,b(\zeta_1)$.
\end{theorem}
\begin{proof}
Suppose $1+\zeta_0\zeta_1=c\,a(\zeta_0)\,b(\zeta_1)$ for all $\zeta$. Evaluating at the four
points $(0,0),(1,1),(0,1),(1,0)$ gives $c\,a(0)b(0)=1$, $c\,a(1)b(1)=2$, $c\,a(0)b(1)=1$, and
$c\,a(1)b(0)=1$. But $\big(c\,a(0)b(0)\big)\big(c\,a(1)b(1)\big)=\big(c\,a(0)b(1)\big)\big(c\,a(1)b(0)\big)$
holds identically, forcing $1\cdot 2 = 1\cdot 1$ --- a contradiction.
\end{proof}
\lstatus{IGL.barrierWitness\_not\_rank\_one}{\proved}

\noindent This is the rigorous core of the informal ``only if'' direction. The fully general
converse (every non-product $\sqrt{|g|}$ blocks factorization) requires non-degeneracy
hypotheses on the function families --- a constant or zero source factorizes anything --- and
is therefore left informal rather than asserted as an axiom.

\begin{theorem}[Compensation restores separability]\label{thm:comp}
Learning the compensated source $\hat f=\tilde f\sqrt{|g|}$ absorbs the metric, so the
integrand is again separable and the Master Formula applies to $\hat f$. The rank cost is
\emph{multiplicative} for metric absorption
($\mathrm{rank}(\tilde f\sqrt{|g|})\le\mathrm{rank}(\tilde f)\,\mathrm{rank}(\sqrt{|g|})$,
Theorem~\ref{thm:rankmul}) and \emph{additive} for the operator correction
($\mathrm{rank}(\hat f)\le\mathrm{rank}(\tilde f\sqrt{|g|})+\mathrm{rank}(L_{\mathrm{coup}}u^*)$,
Theorem~\ref{thm:rankadd}).
\end{theorem}

\section{Cited classical prerequisites}

Mathlib~v4.30.0 contains no theory of Green's functions of elliptic operators and no
exponential-sum approximation theory (verified by source search: no
\leanname{greensFunction} / \leanname{fundamentalSolution} / \leanname{Helmholtz} /
Riesz--Newtonian potential, and no exp-sum results). The two classical facts IGL builds on
are therefore introduced as \emph{named axioms}, with citations, rather than proved. They
are \emph{not} \leanname{sorry}: every theorem using them lists them by name under
\leanname{\#print axioms}, so the dependence on a cited result is always auditable.

\begin{definition}[Green's kernel, cited]\label{ax:green}
For the chosen elliptic operator $L$, the integral kernel $G$ with $u = G * f$ solving
$Lu=f$ is posited per operator.
\lstatus{IGL.greensKernel}{\cited}
\end{definition}

\begin{proposition}[Exponential-sum rank bound, cited]\label{ax:expsum}
There is a constant $C>0$ such that for all $\varepsilon\in(0,1)$ the separable rank
$K(\varepsilon)$ required to approximate an elliptic Green's kernel to accuracy
$\varepsilon$ satisfies $K(\varepsilon) \le C\log(1/\varepsilon) + C$ --- logarithmic in
$1/\varepsilon$ and independent of the dimension (Braess--Hackbusch exponential sums).
\lstatus{IGL.expSumRank\_logBound}{\cited}
\end{proposition}

\section{Formalization status}

\begin{center}
\renewcommand{\arraystretch}{1.25}
\begin{tabular}{@{}llc@{}}
\toprule
\textbf{Result} & \textbf{Lean name} & \textbf{Status}\\
\midrule
Rank-1 factorization (Thm.~\ref{thm:rank1})         & \leanname{igl\_separable\_factorization} & \proved\\
$d$-fold factorization (Thm.~\ref{thm:dfold})       & \leanname{IGL.fubini\_dfold}             & \proved\\
Master Formula (Thm.~\ref{thm:master})              & \leanname{IGL.fubini\_factorization}    & \proved\\
Additive-Green (Cor.~\ref{cor:additive})            & \leanname{IGL.additiveGreen\_factorization} & \proved\\
Kolmogorov--Arnold limit (Cor.~\ref{cor:kan})       & \leanname{IGL.kan\_reduction}           & \proved\\
Gauge invariance (Thm.~\ref{thm:gauge})             & \leanname{IGL.gauge\_invariance}        & \proved\\
Additive eigenvalues $\to$ product (Thm.~\ref{thm:expneg}) & \leanname{IGL.exp\_neg\_sum\_eigs} & \proved\\
Spectral factorization (Thm.~\ref{thm:spectral})    & \leanname{IGL.spectral\_factorization}  & \proved\\
CP-rank sub-additivity (Thm.~\ref{thm:rankadd})     & \leanname{IGL.HasCPRank.add}            & \proved\\
CP-rank sub-multiplicativity (Thm.~\ref{thm:rankmul}) & \leanname{IGL.HasCPRank.mul}          & \proved\\
Rank-friendly differentiation (Thm.~\ref{thm:diffrank}) & \leanname{IGL.applyCoordOp\_hasCPRank} & \proved\\
Barrier, sufficiency (Thm.~\ref{thm:barrierif})     & \leanname{IGL.barrier\_if\_factorization} & \proved\\
Barrier, necessity witness (Thm.~\ref{thm:barrieronly}) & \leanname{IGL.barrierWitness\_not\_rank\_one} & \proved\\
Green's kernel (Def.~\ref{ax:green})                & \leanname{IGL.greensKernel}             & \cited\\
Exp-sum rank bound (Prop.~\ref{ax:expsum})          & \leanname{IGL.expSumRank\_logBound}     & \cited\\
\bottomrule
\end{tabular}
\end{center}

\noindent
All proved results depend only on \leanname{propext}, \leanname{Classical.choice},
\leanname{Quot.sound}; none of them depends on \leanname{IGL.greensKernel} or
\leanname{IGL.expSumRank\_logBound}, confirming that the factorization theory stands on
Mathlib's foundations alone.

\end{document}
